\section{小结}

本报告说明了 Project 3 Roe FVM solver 的方法实现路线。Project 3 保留 Project 2 的整体求解哲学和工程接口：守恒量主推进、原始量同步、CFL 时间步、Riemann case 输入、Tecplot 输出和 exact solver 对比。真正改变的是数值推进核心。

Project 2 使用 MacCormack 预测校正有限差分格式，并通过人工粘性控制间断附近的振荡。Project 3 改用 Godunov 型有限体积格式，其中 Roe approximate Riemann solver 负责在每个界面根据左右状态构造 Roe 平均量、特征值、波强度和右特征向量，得到界面通量 \(\hat{\bm F}_{i+1/2}\)，再通过通量差更新单元平均守恒量。

这种改动使 Project 3 明确成为 Godunov 型 FVM 迎风方法，能够直接体现 Euler 方程的波传播结构。与此同时，Roe 原式也带来新的注意点：它不保证正性，强膨胀和近真空问题可能出现负压或负密度。因此，Project 3 的参数研究不再围绕人工粘性 \(\beta\) 和 sensor 展开，而应围绕 CFL、网格数和 positivity failure 展开。

\section*{参考资料}

\begin{enumerate}[leftmargin=2em]
    \item S. K. Godunov, A difference method for numerical calculation of discontinuous solutions of the equations of hydrodynamics, 1959.
    \item P. L. Roe, Approximate Riemann solvers, parameter vectors, and difference schemes, Journal of Computational Physics, 1981.
    \item E. F. Toro, Riemann Solvers and Numerical Methods for Fluid Dynamics, Springer.
    \item Lesson 11：迎风型格式之 Godunov 与 Roe.
\end{enumerate}
